% % sections/graphical.tex

\subsection*{图文摘要}
% 图文摘要要求：选择1-2幅最能体现论文创新价值的彩图
% 配以英文说明（50个实词以内）

\begin{figure}[H] % 使用 [H] 强制位置
	\centering
	\figbicaption{有限域Montgomery算法工作原理示意图} % 中文标题
	{Schematic diagram of finite field Montgomery algorithm working principle} % 英文标题
	\includegraphics[width=0.6\textwidth,height=0.1\textheight]{figures/figure2.eps}
	% 图片要求：彩色，体现创新价值
\end{figure}

% 英文说明：50个实词以内，突出创新点
Two different Montgomery algorithm implementations show distinct computational pathways. 
The improved version reduces computational steps by 40\% while maintaining accuracy in finite field operations.


% ==================== 编译说明 ====================
% 重要：需要编译两次才能正确显示交叉引用！
% 第一次编译：收集所有标签信息
% 第二次编译：生成正确的引用编号
% 推荐编译顺序：
% 1. xelatex main.tex
% 2. biber main (如果有参考文献)
% 3. xelatex main.tex
% 4. xelatex main.tex

% ==================== 格式要求说明 ====================
% 1. 所有图表标题使用中英文双语
% 2. 表格内容使用英文，表头使用英文
% 3. 图片使用矢量格式（PDF/EPS/AI）
% 4. 采用三线表格式，无竖斜线
% 5. 变量使用斜体，矩阵使用黑斜体
% 6. 图文摘要选择最能体现创新的彩图